% ============================================
% 10. 模型评价与推广
% ============================================

\section{模型评价与推广}

\subsection{模型优点}

\begin{itemize}
  \item \textbf{优点1：模型精度高。}本模型在\textbf{【数据集/场景】}上取得了\textbf{【具体指标数值】}的优异表现，相较于\textbf{【对比方法】}提升了\textbf{【百分比】}，验证了模型的\textbf{【准确性/有效性】}。

  \item \textbf{优点2：模型可解释性强。}通过\textbf{【SHAP分析/特征重要性排序/公式显式表达】}，模型的\textbf{【决策过程/预测依据】}具有清晰的可解释性，便于\textbf{【实际应用/决策支持】}。

  \item \textbf{优点3：模型鲁棒性好。}灵敏度分析结果表明，模型在\textbf{【参数变动范围】}内表现稳定，核心结论不受参数扰动的影响，具有良好的\textbf{【泛化能力/稳健性】}。

  \item \textbf{优点4：建模方法系统完整。}本文从\textbf{【数据处理】}到\textbf{【模型建立】}再到\textbf{【结果验证】}，形成了完整的分析流程，为\textbf{【同类问题】}提供了可参考的范式。
\end{itemize}

\subsection{模型缺点}

\begin{itemize}
  \item \textbf{缺点1：数据依赖性。}模型的性能依赖于\textbf{【数据质量/数据量】}，在\textbf{【数据缺失严重/数据分布变化】}的情况下，模型精度可能\textbf{【下降】}。未来可通过\textbf{【数据增强/迁移学习】}等方法改善。

  \item \textbf{缺点2：计算复杂度。}模型在\textbf{【训练/求解】}过程中需要\textbf{【较多计算资源/较长计算时间】}，对于\textbf{【大规模/实时】}应用场景可能存在\textbf{【效率瓶颈】}。可通过\textbf{【算法优化/并行计算】}加速。
\end{itemize}

\subsection{模型推广}

本文所建立的模型不仅适用于\textbf{【当前赛题场景】}，还可推广至\textbf{【相关领域/更广泛场景】}。例如，在\textbf{【领域A】}中，可将\textbf{【模型核心方法】}应用于\textbf{【具体问题】}；在\textbf{【领域B】}中，通过\textbf{【调整参数/扩展维度】}即可解决\textbf{【类似问题】}。此外，模型的\textbf{【核心思想/技术框架】}对于\textbf{【更广泛的问题类型】}也具有一定的借鉴意义。
